\documentclass[11pt,a4paper]{article}

\usepackage{amsmath,amssymb,mathtools,booktabs,geometry,xurl,enumitem}
\geometry{margin=2.3cm}
\setlist[itemize]{nosep,leftmargin=1.2em}
\renewcommand{\arraystretch}{1.3}

\hyphenation{ap-prox-i-ma-tor ap-prox-i-ma-tion}
\newcommand{\R}{\mathbb{R}}
\newcommand{\N}{\mathbb{N}}
\newcommand{\E}{\mathbb{E}}
\DeclareMathOperator*{\argmin}{arg\,min}

\title{Kernelized polynomial regression\\[0.3em]
\large Formula overview}
\author{Transcribed from Carlos Cotrini,
\emph{Kernelized polynomial regression},\\
Advanced Machine Learning, ETH Zurich, 5 and 8 October 2026}
\date{}

\begin{document}
\maketitle
\thispagestyle{empty}

\noindent
Notation follows the lecture. \(\N=\{0,1,2,\ldots\}\).
A multi-index \(\alpha=(\alpha_1,\ldots,\alpha_d)\in\N^d\) has
\(x^\alpha=\prod_{i\le d}x_i^{\alpha_i}\) and \(\alpha!=\prod_{i\le d}\alpha_i!\).

\tableofcontents
\vspace{1em}

\section{Risk}
\label{sec:risk}

Unknown regression function \(f^*\colon\R^d\to\R\).
Training pairs \((x_i,y_i)\), \(i\le n\), drawn from \(p^*\), with
\begin{equation}
y_i = f^*(x_i)+\varepsilon_i^*,\qquad \varepsilon_i^*\sim\mathcal{N}(0,\sigma^2).
\label{eq:obs}
\end{equation}
The same model is used for a future sample \((x_i^*,y_i^*)\), \(i\le m\).

The quantity to minimize on future data, and its expectation, are
\begin{align}
\hat R_*(f)
&= \frac1m\sum_{i\le m}\bigl(y_i^*-f(x_i^*)\bigr)^2,
\label{eq:future}\\
R(f)
&= \E_{(X,Y)\sim p^*}\bigl[(Y-f(X))^2\bigr].
\label{eq:risk}
\end{align}
Under \eqref{eq:obs} the expected risk splits into approximation and noise,
\begin{equation}
R(f)
= \E\bigl[(f(X)-f^*(X))^2\bigr] + \sigma^2.
\label{eq:risk-split}
\end{equation}
The training substitute is the empirical risk
\begin{equation}
\hat R(f) = \frac1n\sum_{i\le n}\bigl(y_i-f(x_i)\bigr)^2.
\label{eq:empirical}
\end{equation}

\section{Least squares and ridge}
\label{sec:ls}

Two routes give the same linear estimator.
Maximum likelihood with \(y_i\sim\mathcal{N}(\theta^\top x_i,\sigma^2)\) is
\begin{equation}
\hat\theta
= \argmin_{\theta\in\R^d}\sum_{i\le n}\bigl(y_i-\theta^\top x_i\bigr)^2,
\label{eq:mle}
\end{equation}
because the Gaussian likelihood is a decreasing function of that sum.

Statistical learning restricts \(\hat R\) to a hypothesis class and drops the factor \(1/n\), which does not change the minimizer.
The linear class and the ordinary least-squares estimator (OLSE) are
\begin{align}
\mathcal{H}_l
&= \bigl\{f\colon\R^d\to\R \mid f(x)=w^\top x \text{ for some } w\in\R^d\bigr\},
\label{eq:Hl}\\
\hat w
&= \argmin_{w\in\R^d}\sum_{i\le n}\bigl(y_i-w^\top x_i\bigr)^2
= (X^\top X)^{-1}X^\top y,
\label{eq:olse}
\end{align}
with design matrix and response
\begin{equation}
X =
\begin{bmatrix}
x_1^\top \\ \vdots \\ x_n^\top
\end{bmatrix}
\in\R^{n\times d},
\qquad
y =
\begin{bmatrix}
y_1 \\ \vdots \\ y_n
\end{bmatrix}
\in\R^n.
\label{eq:design}
\end{equation}
So \(\hat\theta_{\mathrm{MLE}}=\hat w_{\mathrm{OLSE}}\).
Building \(X^\top X\), building \(X^\top y\), and solving the \(d\times d\) system cost
\(O(nd^2+d^3)\).

If \(y_i=x_i^\top w^*+\varepsilon_i\) with \(\varepsilon_i\sim\mathcal{N}(0,\sigma^2)\), the OLSE is unbiased,
\begin{equation}
\E_{y\mid X}[\hat w] = w^*.
\label{eq:unbiased}
\end{equation}
Among unbiased estimators of the form \(\tilde w=Cy\), it has the smallest variance at every test point (Gauss-Markov):
\begin{equation}
\mathrm{Var}\bigl(x_*^\top\hat w\bigr)
\le \mathrm{Var}\bigl(x_*^\top\tilde w\bigr).
\label{eq:gm}
\end{equation}
Pointwise error still splits as
\begin{equation}
\E\bigl[(y_*-\hat f(x_*))^2\bigr]
= \bigl(\E[\hat f(x_*)]-f^*(x_*)\bigr)^2
+ \mathrm{Var}\bigl(\hat f(x_*)\bigr)
+ \sigma^2.
\label{eq:bv}
\end{equation}
The OLSE term in the middle is often large.
Ridge pays a little bias to shrink it,
\begin{align}
\hat R_r(w)
&= \lVert y-Xw\rVert^2 + \lambda\lVert w\rVert^2,
\label{eq:ridge-obj}\\
\hat w_r
&= \argmin_w \hat R_r(w)
= (X^\top X+\lambda I)^{-1}X^\top y.
\label{eq:ridge}
\end{align}

\section{Polynomial features}
\label{sec:poly}

\(\mathcal{H}_l\) cannot approximate a general continuous target.
On any closed bounded \(B\subset\R^d\), polynomials can: for every continuous \(g\colon B\to\R\) and every \(\delta>0\) there is a polynomial \(f\) with \(\lvert f(x)-g(x)\rvert<\delta\) for all \(x\in B\) \cite{weierstrass1885,stone1948}.
Write \(\mathcal{H}_p\) for that class.

Every \(f\in\mathcal{H}_p\) is linear in a monomial coordinate.
One variable, every degree:
\begin{equation}
f(x)=\sum_{k\in\N} w_k x^k = \langle w,\tilde\varphi(x)\rangle,
\qquad
\tilde\varphi(x)=(1,x,x^2,x^3,\ldots).
\label{eq:phi-1d}
\end{equation}
In \(d\) variables the same statement uses one index per coordinate,
\begin{equation}
f(x)=\sum_{\alpha\in\N^d} w_\alpha x^\alpha = \langle w,\tilde\varphi(x)\rangle.
\label{eq:phi-d}
\end{equation}
Raw monomials are unusable as features: at \(x=1\) and \(w=(1,1,\ldots)\) the inner product diverges, and for \(d=100\) the monomials of degree \(5\) already number almost \(10^8\).

Scale each monomial by a Gaussian bump and a factorial,
\begin{equation}
\varphi(x)
= e^{-\lVert x\rVert^2/2}
\Biggl[\frac{x^\alpha}{\sqrt{\alpha!}}\Biggr]_{\alpha\in\N^d}.
\label{eq:phi}
\end{equation}
Then \(\lVert\varphi(x)\rVert=1\) for every \(x\), there is no degree cutoff, and the model remains a universal approximator.
The fit is still linear regression on features,
\begin{equation}
f(x)=\langle w,\varphi(x)\rangle.
\label{eq:poly-fit}
\end{equation}
The one-variable map is the same formula with \(\lVert x\rVert^2=x^2\) and \(\alpha=k\).

\section{The kernel trick}
\label{sec:kernel}

Stack the feature rows \(\Phi\in\R^{n\times\infty}\), with row \(i\) equal to \(\varphi(x_i)^\top\).
Polynomial least squares asks for
\begin{equation}
\hat w_\infty
= \argmin_{w\in\R^\infty}\lVert\Phi w-y\rVert^2.
\label{eq:poly-ls}
\end{equation}
The normal equation \((\Phi^\top\Phi)^{-1}\Phi^\top y\) is an infinite matrix.
The minimum-norm solution never forms it,
\begin{equation}
\hat w_\infty = \Phi^\top(\Phi\Phi^\top)^{-1} y.
\label{eq:min-norm}
\end{equation}
(The lecture treats the usual normal equation as shorthand: \(\Phi^\top\Phi\) has no inverse. Equation~\eqref{eq:min-norm} is the exact minimum-norm solution.)

Prediction at \(x_*\in\R^d\) uses only \(n\times n\) kernel entries.
The sum over all scaled monomials collapses to the Gaussian kernel,
\begin{equation}
k(x,x')
= \langle\varphi(x),\varphi(x')\rangle
= \exp\Bigl(-\frac12\lVert x-x'\rVert^2\Bigr).
\label{eq:gauss}
\end{equation}
With \(K_{ij}=k(x_i,x_j)\) and \(k_{*,i}=k(x_*,x_i)\),
\begin{equation}
\hat f(x_*)
= \varphi(x_*)^\top\Phi^\top(\Phi\Phi^\top)^{-1} y
= k_*^\top K^{-1} y.
\label{eq:k-pred}
\end{equation}
The same identity with a ridge penalty is exact, not a minimum-norm convention,
\begin{align}
\hat w_\infty
&= \argmin_{w\in\R^\infty}
\Bigl(\lVert\Phi w-y\rVert^2+\lambda\lVert w\rVert^2\Bigr)
= \Phi^\top(\Phi\Phi^\top+\lambda I)^{-1} y,
\label{eq:k-ridge-w}\\
\hat f(x_*)
&= k_*^\top (K+\lambda I)^{-1} y.
\label{eq:k-ridge}
\end{align}
Algorithm: build \(K\) from the training points, solve \((K+\lambda I)\beta=y\), then for each test point build \(k_*\) and return \(k_*^\top\beta\).
Cost \(O(n^2 d+n^3)\). At \(n=10^6\) that is about \(10^{18}\) operations and on the order of \(8\,\mathrm{TB}\) for \(K\).

\section{Random features}
\label{sec:rf}

Keep the predictor, replace the infinite map by \(2Q\) numbers with \(2Q<n\).
Draw frequencies \(\omega_q\sim\mathcal{N}(0,I_d)\) once, before seeing data, and set
\begin{equation}
\psi(x)
= \frac1{\sqrt Q}
\begin{pmatrix}
\cos(\omega_1^\top x) \\
\sin(\omega_1^\top x) \\
\vdots \\
\cos(\omega_Q^\top x) \\
\sin(\omega_Q^\top x)
\end{pmatrix}
\in\R^{2Q}.
\label{eq:psi}
\end{equation}
Rows \(\psi(x_i)^\top\) form \(\Psi\in\R^{n\times 2Q}\). The estimator is ordinary ridge regression in these coordinates,
\begin{equation}
\hat f(x_*)
= \psi(x_*)^\top(\Psi^\top\Psi+\lambda I)^{-1}\Psi^\top y.
\label{eq:rf}
\end{equation}
Cost \(O(nQd+nQ^2+Q^3)\), linear in the sample size.

It matches kernel ridge when \(\psi(x)^\top\psi(x')\) matches \(k(x,x')\).
Bochner's theorem writes the Gaussian kernel as a Fourier expectation \cite{bochner1933,rudin1962}.
Rahimi and Recht turn that expectation into features \cite{rahimi2007}.
Because \(k\) is real and even, the expectation is an average of cosines,
\begin{equation}
k(x,x')
= \kappa(x-x')
= \E_{\omega\sim\mathcal{N}(0,I_d)}
\bigl[\cos\bigl(\omega^\top(x-x')\bigr)\bigr].
\label{eq:bochner}
\end{equation}
The average of \(Q\) independent draws is
\begin{equation}
k(x,x')
\approx \frac1Q\sum_{q\le Q}\cos\bigl(\omega_q^\top(x-x')\bigr)
= \psi(x)^\top\psi(x').
\label{eq:lln}
\end{equation}
The cosine splits by the angle formula, which is why \(\psi\) stores both \(\cos(\omega_q^\top x)\) and \(\sin(\omega_q^\top x)\).
The weak law puts the error near \(1/\sqrt{2Q}\). At \(Q=5000\) that is about \(0.01\).

Then \(K\approx\Psi\Psi^\top\) and \(k_*\approx\Psi\psi(x_*)\).
The push-through identity
\begin{equation}
\Psi^\top(\Psi\Psi^\top+\lambda I)^{-1}
= (\Psi^\top\Psi+\lambda I)^{-1}\Psi^\top
\label{eq:push}
\end{equation}
turns the \(n\times n\) kernel predictor into the \(2Q\times 2Q\) estimator \eqref{eq:rf}.

\begin{table}[ht]
\centering
\caption{Three estimators of the same squared loss.}
\label{tab:costs}
\begin{tabular}{lll}
\toprule
Estimator & Coordinates & Work \\
\midrule
Linear ridge & \(x\in\R^d\) & \(O(nd^2+d^3)\) \\
Kernel ridge & \(\varphi(x)\in\R^\infty\) & \(O(n^2 d+n^3)\) \\
Random features & \(\psi(x)\in\R^{2Q}\) & \(O(nQd+nQ^2+Q^3)\) \\
\bottomrule
\end{tabular}
\end{table}

\noindent
The kernel buys every scaled monomial. Random features buy the same predictor back at a cost linear in \(n\).

\section{Source index}
\label{sec:sources}

Formulas in this note are transcribed from the lecture slides.
Named theorems are the ones the slides cite.
Paths below are files in this repository.

\begin{table}[ht]
\centering
\caption{Internal sources.}
\label{tab:sources}
\begin{tabular}{ll}
\toprule
Role & Path \\
\midrule
Lecture slides & \path{experiments/006/kernelized-polynomial-regression-slides.pdf} \\
This note & \path{experiments/006/formulas/main.tex} \\
\bottomrule
\end{tabular}
\end{table}


{\small
\bibliographystyle{plain}
\bibliography{references}
}

\end{document}
